% Generated by scripts/tikz_reviews.py using the compiled book's labels.
\pdfbookmark[0]{Chapter 11}{comparison-chapter-11}
\ReviewFigure{fig-polytope-proof}
{Interior points and failure of exact recovery}
{figures/sparse-theory/proofs/polytope-proof-1.png}
{figures/tikz/sparse-theory-polytope-proof.pdf}
{The onto-map polytope characterization: x0 solves basis pursuit iff Ax0/||x0||1 lies on the boundary of AB.}
{Uses the normalized polytope AB and q=Ax0/r, r=||x0||1, consistently with the corrected proof. The radial extension (1+t)q is constructed exactly on the ray through q and on the polygon boundary. The second panel shows h as a directed displacement to q+h inside an interior ball. Both panels illustrate failure of optimality through interior membership; no uniqueness claim is implied.}
{sparse-theory-polytope-proof}
{Complete figure}
{11.2}
\ReviewFigure{fig-review-sparse-theory-injectivity}
{Removing a dependent active coefficient}
{figures/sparse-theory/proofs/injectivity.png}
{figures/tikz/sparse-theory-injectivity.pdf}
{Support reduction along x(t)=x+t h when A\_I h\_I=0.}
{Three affine coefficient trajectories have a marked interval with both endpoints shown. The first coefficient vanishes at the left endpoint and the third at the right. The figure states that x is a minimizer: this assumption forces the affine l1 norm to have zero slope while signs stay fixed. Constant fidelity follows from Ah=0. The norm remains constant at the endpoints by continuity, though the sign pattern changes there.}
{sparse-theory-injectivity}
{Complete figure}
{11.4}
\ReviewFigure{fig-bregman}
{Bregman divergence as a vertical gap}
{figures/sparse-theory/proofs/bregman.png}
{figures/tikz/sparse-theory-bregman.pdf}
{Definition of D\_eta(x|x0) for eta in the subdifferential of a convex J.}
{The left panel uses a strictly convex quadratic and its tangent; the positive vertical gap is the Bregman divergence. The right panel uses J=abs(x), with x0 and x on the same positive affine branch, so the divergence vanishes even when x differs from x0.}
{sparse-theory-bregman}
{Complete figure}
{11.5}
\ReviewFigure{fig-review-sparse-theory-bregman-l1}
{A strict subgradient margin controls the error}
{figures/sparse-theory/proofs/bregman-l1.png}
{figures/tikz/sparse-theory-bregman-l1.pdf}
{The coordinatewise estimate in the l1 Bregman lower bound.}
{Explicitly sets x0=0, as required on nonsaturated coordinates for J=abs(x). For abs(eta)<1 the supporting line eta*x lies strictly below abs(x) away from zero. The lower bound is D\_eta(x|0)>= (1-abs(eta))*abs(x), not a norm bound on saturated coordinates.}
{sparse-theory-bregman-l1}
{Complete figure}
{11.6}
\ReviewFigure{fig-review-sparse-theory-recovery-region}
{Parameter region for sign consistency}
{figures/sparse-theory/proofs/recov-zone.png}
{figures/tikz/sparse-theory-recovery-region.pdf}
{Intersection of delta+lambda<T/K and R*delta<S*lambda in the support-stability proof.}
{Reconstructs the two strict inequalities and the contained simpler triangle. The vertical cutoff is lambda\_max=TR/[K(R+S)], using R=KL+1. The oblique boundaries and the cutoff for the smaller triangle are excluded; the intersection marker is open. The zero-noise axis delta=0 is allowed for positive lambda satisfying the displayed strict inequalities. The horizontal axis is lambda and the vertical axis is delta.}
{sparse-theory-recovery-region}
{Complete figure}
{11.8}
\ReviewFigure{fig-spikes}
{Convolution of a discrete measure}
{figures/sparse-theory/proofs/spikes.png}
{figures/tikz/sparse-theory-convolution-spikes.pdf}
{The dictionary map phi(z), discrete measure m\_x, and measurement operator A.}
{Two translated Gaussian kernels have unit peak; the spike heights and the output use the same vertical scale, so the displayed output is exactly their weighted sum. Kernel centers, spike locations, and output-location labels agree across the panels. The diagram distinguishes index-space weights from physical reconstruction-grid locations, and writes the output as both the continuous operator applied to the discrete measure and the dictionary product Ax.}
{sparse-theory-convolution-spikes}
{Complete figure}
{11.10}
